\noindent {\sffamily\small\bfseries\color{C3}PART II · EXPERIMENT DECISION PROCESS}
\section{先把“比较什么”说清楚}
\subsection{材料审读先于正式计算}
Codex首先检查了已有仓库和两个新ZIP，保留原件与实验一成果。中文文件名存在混合误解码，安全提取时记录名称映射、内容哈希和无意义元数据过滤；两城市数据在多个归档里各有三份，核对后确认是重复副本，不能增加统计样本数。两份PPT共39页、两份Notebook共59 Cells分别完整审读，教师九道进阶思考题及基础讨论逐项登记。

审读暴露了两个不能靠完成代码掩盖的问题。基础Notebook有实际TODO5，但文字只说TODO1--4；BB的Assignment代码已经存在，真正的风险是把未匹配队列清空。课件还提瓶颈目标，Notebook却用总成本KM。Codex没有任选一种解释覆盖另一种，而是分别保留sum和max实现及参考。这是当前提示词已明确要求的核验，不是后来为报告补造的争论。

\subsection{低延迟曲线为何先看覆盖}
首轮55个输入、1045条方法指标中，教师场景短批次discard只服务25\%请求，司机先到的场景甚至没有服务；回收版本则能完成可用的全部供给。这个读数使图型选择有了明确目的：如果只画已服务订单的分布，丢单可能看起来更好。

Codex实际生成了完整分母ECDF和条件raincloud两种候选。前者让未服务质量保留在较低平台，后者虽能显示分布却必须另印覆盖分母。因此正式主图选择ECDF加覆盖/尾部对照，而不是因为某种图更复杂。统计单位仍是每个实例及其请求，不能把同seed内订单当作更多独立实验。

\subsection{从兜底比例到阈值诊断}
小阈值0.5/1/2有平均17\%匹配留到horizon，阈值5为7\%。Codex据此提出机制假设：过小阈值阻塞早期匹配，过大阈值又容易成为局部即时选择。随后才在3种到达顺序、相同5seed上细化阈值2/3/4/5/6/8/10/15，没有改变TBG核心或凭空增加deadline事件。

另一条负例也被保留：一个批次内最小最大成本的BBmax，在teacher seed13全天最大延迟比BBsum高2.162 min。独立核验确认实现与参考正确，因此它是未来司机选择带来的在线局限，不需要把算法改到每个指标都赢。最终独立合成seed110--114只用于确认，不重新回流选参。

\section{进化优化怎样从端点不足走到有限组合}
\subsection{先比较基础预算，再引入初始化}
进化分支先验证小图编码、解码、非支配关系与二维/三维HV，然后运行教师100种群/300代/5seed基线。18点开发批次中，NSGA-II的HV为1.090244，NSGA-III为1.077092，加权KM为1.098977。这些读数与课件参考数值分开保存，高收入端差距触发下一轮，而不是预先决定复杂混合策略更好。

先增加种群或代数，只改变一个因素。NSGA-II增种群比增代数更有利，而NSGA-III相反，说明同样约六万次评价的资源增加并非相同干预。之后才把真实KM匹配编码成初始个体，保持原100/300预算，避免把收益混在更多评价次数里。

KM单项将II/III的HV提高至1.111427/1.111494，并减小seed波动。接着测试它与种群200的一个有限组合，保留基线、增种群、KM单项和组合四个父项。组合只比KM单项增加很小HV，却增加约一倍EA评价、Spacing变差，因此没有升级最终确认预算。五个未来代理批次上只确认冻结的II原预算与KM单项，五对HV差均为正，均差0.031470。这是所测输入的支持，不是任意未来需求的理论保证。

\subsection{独立审核修正评价，不掩盖旧结果}
Research核验者发现，KM全部离散权重输出中有支配点，约$10^{-13}$求和差还会制造几乎重复的点。原始输出保持不变，报告指标统一规范去重再取非支配集，并做8/10/12位精度检查。此前Spacing不再作为当前有效评价，但旧运行没有被改写成从未发生。

初始化额外41次标量求解的成本则在正式hybrid运行前补入。这个修复影响公平预算解释，不是新增科研增益。NSGA-III在球面DTLZ2上更合适，在三目标派单细带上的优势不明显，也没有为迎合“目标更多就更好”的故事而删去。
